% Source: https://artofproblemsolving.com/community/c3h2360247_math_advent_calendar
% Stable identity: aops-topic-2360247-q04
The centres are $O_1O_2=3+3=6$ apart, and $T$ is the midpoint of $O_1O_2$. The large circle with centre $T$ and radius $R$ is internally tangent to the circle with centre $O_1$ and radius $3$ exactly when $R-3=TO_1=3$, so $R=6$.

\begin{asy}
usepackage("amsmath");
size(4.5cm);
defaultpen(fontsize(10pt));
pair T=(0,0), O1=(-3,0), O2=(3,0), P=(0,6), Q=(0,-6);
draw(circle(O1,3)); draw(circle(O2,3)); draw(circle(T,6),linewidth(0.9));
draw(O1--P--O2--Q--cycle,blue+linewidth(0.9));
draw(O1--O2,dashed); draw(P--Q,dashed);
dot(T); dot(O1); dot(O2); dot(P); dot(Q);
label("$T$",T,dir(-45)); label("$O_1$",O1,W); label("$O_2$",O2,E); label("$P$",P,N); label("$Q$",Q,S);
label("$3$",(-1.5,0),N); label("$6$",(0,3),E);
\end{asy}

The diagonals of a rhombus are perpendicular and bisect each other. One diagonal is $O_1O_2$, with midpoint $T$; so the other diagonal $PQ$ lies on the line through $T$ perpendicular to $O_1O_2$, with $T$ as its midpoint. Since $P$ and $Q$ lie on the large circle, $TP=TQ=6$ and $PQ=12$. The area of a rhombus is half the product of its diagonals:
\[\dfrac12\cdot6\cdot12=\boxed{36}.\]
